% Overleaf-Datei (Name zu bestaetigen), Spiegel in docs/overleaf/.
% Nur zusammen mit der Overleaf-Datei aendern.
% Inhalt: Abstract.

\section*{Abstract}
\addcontentsline{toc}{section}{Abstract}

Machine-learned interatomic potentials for reactions are trained on DFT
labels, and the level of theory of those labels bounds their accuracy.
The common way to raise the level of theory is relabelling: the geometries of an
existing dataset are kept and only their energies and forces are
recomputed at a higher level. This thesis asks two questions about
that practice. RQ1: can a model be lifted to a higher level by
relabelling only a subset of its training data? RQ2: can a model learn
a surface from paths that were never relaxed on that surface?

For RQ1, a MACE model is trained on Transition1x
(\mbox{$\omega$B97X/6-31G(d)}), and a correction head on top of it
is trained to predict the difference to \mbox{$\omega$B97M-V/def2-TZVP}
from about 80\,000 geometries, under 1\,\% of the dataset, relabelled
restricted at that level. It is tested on 30 reactions grouped by low, mid and high
multireference character, at fixed geometries and in transition-state
searches the model drives itself. Forces improve everywhere, and the
transition-state searches succeed on every reaction. Barriers improve
threefold for reactions of low multireference character; for reactions
of mid and high multireference character the correction overshoots,
turning the overestimation into an underestimation of similar size,
and at high character the corrected barriers are worse than the
uncorrected ones. The answer to RQ1 is yes, with the reservation of
the reactions of high multireference character.

In RQ1 the labels were raised to a higher level of theory while the
spin formalism stayed restricted. In RQ2 we ask whether relabelling still
works when the spin formalism changes as well. For that we turn to OMol25, the
largest molecular DFT dataset to date with more than 100 million
structures labelled unrestricted at \mbox{$\omega$B97M-V/def2-TZVPD}.
No reaction path in OMol25 was relaxed on the surface the labels were
computed on. This includes the Transition1x paths, which were taken
over unchanged from restricted searches at a lower level. Published
benchmarks of the OMol25 models report high success rates in
transition-state search. None has tested the models at
broken-symmetry transition states. This work closes that gap for three
of the models: UMA-S, UMA-M and eSEN. We let them search the
transition states of 45 Transition1x reactions. Each transition state
the models deliver is evaluated with an unrestricted
DFT single point and a stability analysis, which sorts the transition
states into closed-shell and broken-symmetry. The models predict single-point energies well for both groups. But
their force errors are two to three times larger at the
broken-symmetry transition states. The structures they deliver there
are therefore not stationary points of the unrestricted surface, with
residual forces up to 3.4 times those of the closed-shell group, and
the models fail to predict the barriers of these reactions
accurately. The training data explain this: we measure that the training points
of these reactions lie away from the transition state of the
unrestricted surface, so the models have seen the region around the
broken-symmetry transition state and never the transition state
itself.

Two things follow for practice: a dataset that relabels paths with a
different spin formalism should re-optimise them on the new surface, and a stability
analysis should precede any further use of a transition state whose
multireference character cannot be ruled out.
